\bookchapter{الفصل 45: ضمان الجودة — اجعل كود الذكاء الاصطناعي موثوقا}

\chapterintro{لكود الذكاء الاصطناعي سمة لافتة: فهو “يبدو صحيحا” في الغالب — أسماء متغيرات منظمة وتعليقات كاملة ومنطق واضح. لكن “أن يبدو صحيحا” لا يعني “أن يكون صحيحا فعلا”. المشاكل الخفية في كود الذكاء الاصطناعي غالبا أصعب اكتشافا من الكود المكتوب يدويا، لأن جودته السطحية عالية جدا فتدعو إلى التهاون. يعلّمك هذا الفصل كيف تبني نظاما فعالا لضمان جودة كود الذكاء الاصطناعي — بثلاثة خطوط دفاع هي الاختبارات والمراجعة والتحكم في الإصدارات، لتتحول “يبدو صحيحا” إلى “صحيح فعلا”.}

\booksection{أولاً: التوجيه بالاختبارات — دع الذكاء الاصطناعي يكتب الاختبارات}

الاختبار هو خط الدفاع الأول لجودة البرمجة بالذكاء الاصطناعي. والخبر السار: جعل الذكاء الاصطناعي يكتب الاختبارات أسهل وأوثق من جعله يكتب كود الوظائف — لأن السلوك المتوقع في الاختبارات محدد بوضوح، ومن النادر أن ينتج كودا “يبدو صحيحا وهو ليس كذلك”.

سير العمل الموجه بالاختبارات في البرمجة بالذكاء الاصطناعي هو: اطلب أولا من الذكاء الاصطناعي كتابة كود الوظائف، وبعد نجاح التحقق الأساسي اطلب منه فورا كتابة حالات الاختبار. ينبغي أن تغطي حالات الاختبار المسار الطبيعي والحالات الحدودية والحالات الاستثنائية. ثم اطلب من الذكاء الاصطناعي تصحيح كود الوظائف بناء على نتائج الاختبارات. ميزة هذا السير أن الاختبارات بمجرد كتابتها تصبح أداة تحقق دائمة — ويمكن في أي تعديل لاحق التأكد من عدم إحداث مشاكل انحدار (Regression) بتشغيل الاختبارات.

\begin{lstlisting}
اكتب اختبارات وحدة للدالة calculate_discount.
يجب أن تغطي حالات الاختبار:
1. المسار الطبيعي:
   - 100 يوان بخصم 20% -> 80
   - 200 يوان بخصم 50% -> 100
2. الحالات الحدية:
   - 0 يوان بخصم 20% -> 0
   - 100 يوان بدون خصم -> 100 (لا تخفيض)
   - 100 يوان بالسعر الكامل -> 100 (السعر الكامل)
3. حالات استثنائية:
   - سعر سالب -> يجب أن يرمي خطأً
   - نسبة الخصم > 1 -> يجب أن يرمي خطأً
   - إدخال غير رقمي -> يجب أن يرمي خطأً
استخدم إطار Vitest وأعطني ملف الاختبارات الكامل.
\end{lstlisting}

لاحظ خصائص هذه التعليمات: تحدد صراحة ثلاث فئات من حالات الاختبار (طبيعية وحدودية واستثنائية)، ولا تترك للذكاء الاصطناعي حرية “الاختيار الانتقائي” لما يختبره. قد يكتب الذكاء الاصطناعي أحيانا اختبارات المسار الطبيعي وحدها “لإنجاز المهمة بسرعة”، وتحديد كل الفئات صراحة يجنبنا هذه المشكلة.

\booksection{ثانياً: مهارات مراجعة الكود}

حتى مع وجود تغطية اختبارات، يبقى عليك مراجعة كود الذكاء الاصطناعي بنفسك. مراجعة الكود ليست قراءة سطرا بسطر — فهذا مكلف زمنيا ويعتمد كثيرا على مستواك البرمجي. المراجعة الفعالة لكود الذكاء الاصطناعي تركز على الأبعاد الأساسية الآتية:

\begin{center}
\begin{tabularx}{\linewidth}{|X|X|X|}
\hline
\textbf{بُعد المراجعة} & \textbf{ماذا نراقب} & \textbf{المشاكل الشائعة} \tabularnewline
\hline
الأمان & التحقق من المدخلات، معالجة المفاتيح، الحماية من حقن SQL & غالبا ما ينسى الذكاء الاصطناعي التحقق من المدخلات، ويكتب المفاتيح داخل الكود مباشرة \tabularnewline
معالجة الأخطاء & هل تلتقط الاستثناءات؟ وهل رسائل الخطأ معقولة؟ & يميل الذكاء الاصطناعي إلى كتل catch فارغة أو console.log فقط \tabularnewline
الأداء & هل توجد مشاكل أداء واضحة مثل استعلامات N+1 أو حلقات لا نهائية & أحيانا يكتب الذكاء الاصطناعي خوارزميات بتعقيد O(n²) \tabularnewline
قابلية الصيانة & هل هيكل الكود واضح؟ وهل الأسماء ذات معنى؟ & أحيانا يولد الذكاء الاصطناعي كودا مُجردا بإفراط \tabularnewline
إدارة الاعتماديات & هل الاعتماديات المستوردة ضرورية؟ وهل الإصدارات محددة بوضوح؟ & يميل الذكاء الاصطناعي إلى استيراد مكتبات غير ضرورية \tabularnewline
\hline
\end{tabularx}
\end{center}

\begin{center}\small\color{gray}الأبعاد الأساسية لمراجعة كود الذكاء الاصطناعي\end{center}

لا داعي لفحص كل شيء بنفسك أثناء المراجعة — يمكنك تكليف الذكاء الاصطناعي بالجولة الأولى من المراجعة. بتعليمة مثل: “راجع أمان هذا الكود ومعالجة أخطائه وأداءه، وأشر إلى المشاكل المحتملة”. تتقدم قدرة الذكاء الاصطناعي على المراجعة الذاتية باستمرار، ورغم أنها لا تغني كليا عن المراجعة البشرية، فإنها تساعدك على تصفية معظم المشاكل الواضحة.

\booksection{ثالثاً: التحكم في الإصدارات — دور Git في البرمجة بالذكاء الاصطناعي}

دور Git في البرمجة بالذكاء الاصطناعي أهم منه في البرمجة التقليدية. فبما أن كل تعديل يجريه الذكاء الاصطناعي قد يقدم تغييرات غير متوقعة، تحتاج إلى “دواء الندم” الموثوق — والتحكم في الإصدارات هو ذلك الدواء.

أفضل ممارسة لـ Git في البرمجة بالذكاء الاصطناعي هي “إرسال متكرر وتقدم بخطوات صغيرة”. أرسل Commit بعد إنجاز كل ميزة قابلة للتحقق. هكذا إذا سبب أحد تعديلات الذكاء الاصطناعي مشكلة، تستطيع التراجع بدقة إلى آخر حالة سليمة دون هدم كل عملك.

وحيلة عملية أخرى هي “العزل بالفروع”. للتعديلات التي لا يمكن الجزم بنتائجها (مثل إعادة الهيكلة الكبيرة)، أنشئ فرعا جديدا للعمل عليها أولا. إن أعجبتك النتيجة ادمجت الفرع في الفرع الرئيسي؛ وإلا حذفت الفرع مباشرة دون المساس بالرئيسي. هذا أكثر أمانا بكثير من التعديل المتكرر على الفرع الرئيسي.

\begin{enumerate}
  \item قبل كل تعديل يجريه الذكاء الاصطناعي، تأكد من نظافة مساحة العمل (لا تغييرات غير مُرسلة). هكذا إذا لم تعجبك تعديلات الذكاء الاصطناعي، يمكنك التراجع بسرعة بـ git checkout .
  \item بعد كل تعديل من الذكاء الاصطناعي ونجاح التحقق، أرسل Commit فورا. وليكن نص الإرسال يصف ما أجراه هذا التعديل، مثل "feat: إضافة وظيفة تسجيل المستخدم".
  \item لإعادة الهيكلة الكبيرة، أنشئ فرعا أولا: git checkout -b refactor/auth-system. أنجز التعديلات على الفرع، وبعد نجاح الاختبارات ادمجه.
  \item اكتسب عادة مراجعة الـ diff. بعد كل تعديل من الذكاء الاصطناعي، اعرض ما تغير تحديدا بـ git diff. هذا لا يساعدك على اكتشاف المشاكل فحسب، بل يساعدك أيضا على تعلم الكود.
\end{enumerate}

\booksection{رابعاً: الفخاخ الشائعة ودليل تجنبها}

في البرمجة بالذكاء الاصطناعي فخاخ تتكرر باستمرار ويقع فيها جميع المبتدئين تقريبا. معرفتها مسبقا تجنبك الكثير من الالتفافات الخاطئة.

\booksubsection{الفخ الأول: فخ الثقة العمياء}

“كود الذكاء الاصطناعي يبدو احترافيا، فلا بد أنه سليم؟” — هذا أخطر الفخاخ. الكود الذي يولده الذكاء الاصطناعي نظيف جدا ظاهريا — أسماء منظمة وتعليقات كاملة وبنية واضحة — وهذه “المظهرية الاحترافية” تدعو إلى التهاون. لكن النظافة الظاهرية لا تعني سلامة المنطق. الترياق: تحقق دائما بالاختبارات، ولا تحكم بمجرد “أنه يبدو صحيحا”.

\booksubsection{الفخ الثاني: فخ التضخم}

“ما دام الذكاء الاصطناعي يستطيع كتابة الكود لي، فلماذا لا نجعل الوظائف شاملة أكثر؟” — هذا يؤدي إلى تضخم الكود. الذكاء الاصطناعي بارع في “الإفراط في العطاء”: تطلب منه صفحة تسجيل دخول فيمنحك نظام إدارة مستخدمين كاملا. الوظائف الزائدة لا تزيد التعقيد فحسب، بل ترفع احتمال ظهور الأخطاء. الترياق: التزم بصرامة بمبدأ Simplicity First، وأخبر الذكاء الاصطناعي صراحة “نفذ هذه الوظيفة فقط ولا تضف شيئا إضافيا”.

\booksubsection{الفخ الثالث: فخ الفهم}

“الذكاء الاصطناعي المفترض أن يفهم قصدي؟” — ليس بالضرورة. رغم قوة فهم الذكاء الاصطناعي، فإنه “يخمن” حين تكون المتطلبات غامضة. وإذا أخطأ التخمين كان الكود الناتج خاطئا. والأعسر أن كوده الخاطئ “يبدو معقولا” عادة، فيوهمك أنه فهم. الترياق: قدم في التعليمات تفاصيل محددة كافية لتقليل المجال الذي يضطر فيه الذكاء الاصطناعي إلى “التخمين”.

\booksubsection{الفخ الرابع: فخ السياق}

“قلت ذلك في حوار سابق، فالذكاء الاصطناعي ما زال يتذكر؟” — من المحتمل جدا أنه لا يتذكر. مع طول الحوار أو فتح جلسة جديدة، يفقد الذكاء الاصطناعي السياق السابق. وإذا لم يكتب قيد أساسي في ملف ذاكرة المشروع، فمن المرجح جدا إغفاله في التعديلات اللاحقة. الترياق: اكتب القيود الأساسية في ملف ذاكرة المشروع، ولا تعتمد على ذاكرة الحوار.

\begin{tcbwarn}{تذكير سريع بالفخاخ الأربعة}
فخ الثقة — لا تصدق “أنه يبدو صحيحا”؛ فخ التضخم — الذكاء الاصطناعي يحب “الإفراط في العطاء”؛ فخ الفهم — مع المتطلبات الغامضة سـ“يخمن”؛ فخ السياق — ذاكرة الحوار سـ“تنسى”. والترياقات المقابلة: التحقق بالاختبارات، وقضم الزائد، وتقديم تفاصيل كافية، والكتابة في ملف الذاكرة.

\end{tcbwarn}

\begin{tcbpractice}{لنطبق عملياً: ابن خطوط الدفاع الثلاثة لمشروعك}
إذا كنت قد بدأت المشروع الأول (موقع الأدوات): (1) اطلب من الذكاء الاصطناعي كتابة 3 فئات من حالات الاختبار للوظائف الأساسية؛ (2) راجع الكود وفق الأبعاد الخمسة (الأمان/معالجة الأخطاء/الأداء/قابلية الصيانة/الاعتماديات)؛ (3) هئ مشروعك بـ git وأجر أول Commit. حين تكتمل خطوط الدفاع الثلاثة يصبح مشروعك ملكك حقا.

\end{tcbpractice}

\begin{tcbsummary}{خلاصة الفصل}
\begin{itemize}
  \item كود الذكاء الاصطناعي "يبدو صحيحا" لا يعني "صحيح فعلا" — فجودته السطحية العالية هي ما يدعو غالبا إلى التهاون
  \item دع الذكاء الاصطناعي يكتب الاختبارات: حدد صراحة الفئات الثلاث (طبيعية/حدودية/استثنائية) لمنع "الاختيار الانتقائي"
  \item أبعاد المراجعة الخمسة: الأمان، معالجة الأخطاء، الأداء، قابلية الصيانة، إدارة الاعتماديات
  \item Git هو "دواء الندم": إرسال متكرر، وتقدم بخطوات صغيرة، وعزل بالفروع، ومراجعة الـ diff باستمرار
  \item الفخاخ الأربعة: الثقة (التحقق بالاختبارات)، التضخم (قضم الزائد)، الفهم (تفاصيل كافية)، السياق (الكتابة في ملف الذاكرة)
\end{itemize}

\end{tcbsummary}
